% Generated by `taskgen course examples` from ecc.point_order (seed 5).
% Do not edit: change examples.yaml and regenerate.
\begin{example}[порядок точки]\label{ex:ec-point-order}
Найти порядок точки $P = (6, 16)$ эллиптической кривой \[E\colon\ y^2 = x^3 + 21x + 1\quad\text{над } \mathbb{F}_{29},\qquad \#E(\mathbb{F}_{29}) = 39,\] последовательно вычисляя кратные $kP$. Согласуется ли ответ с теоремой Лагранжа?

\par\noindent\textit{Решение.} Если $\ord P = n$, то $(n - k)P = -kP$. Поэтому достаточно вычислять $kP = (k - 1)P + P$, пока не встретится точка с $y = 0$ (тогда $2kP = \mathcal{O}$ и $n = 2k$) или $x(kP) = x((k-1)P)$, т.\,е. $kP = -(k-1)P$ (тогда $n = 2k - 1$): раньше этого ни одно из равенств $n \mid 2k$, $n \mid 2k - 1$ выполняться не может. Всё по модулю $29$, $a = 21$.\[\begin{array}{r|l|r|c}k & \lambda & \lambda \bmod 29 & kP \\ \hline 1 & & & (6, 16) \\ 2 & \frac{3 \cdot 6^2 + 21}{2 \cdot 16} \equiv \frac{13}{3} & 14 & (10, 15) \\ 3 & \frac{16 - 15}{6 - 10} \equiv \frac{1}{25} & 7 & (4, 27) \\ 4 & \frac{16 - 27}{6 - 4} \equiv \frac{18}{2} & 9 & (13, 8) \\ 5 & \frac{16 - 8}{6 - 13} \equiv \frac{8}{22} & 3 & (19, 3) \\ 6 & \frac{16 - 3}{6 - 19} \equiv \frac{13}{16} & 28 & (5, 12) \\ 7 & \frac{16 - 12}{6 - 5} \equiv \frac{4}{1} & 4 & (5, 17) \\ \end{array}\](Строка $k = 2$ --- удвоение, $\lambda = \frac{3x^2 + a}{2y}$; остальные --- сложение с $P$, $\lambda = \frac{y_2 - y_1}{x_2 - x_1}$. Обратные: $3^{-1} \equiv 10$, $25^{-1} \equiv 7$, $2^{-1} \equiv 15$, $22^{-1} \equiv 4$, $16^{-1} \equiv 20$, $1^{-1} \equiv 1$.)\par $x(7P) = x(6P)$, т.\,е. $7P = -6P$ и $13P = \mathcal{O}$: $\ord P = 2 \cdot 7 - 1 = 13$.\par Теорема Лагранжа: порядок точки делит порядок группы, $39 = 13 \cdot 3$ --- согласуется.\[\ord P = 13.\]
\end{example}
